\documentclass[11pt]{article}

\usepackage[
  top    = 1.1in,
  bottom = 1.1in,
  left   = 1.25in,
  right  = 1.25in
]{geometry}

\usepackage{microtype}
\usepackage{lmodern}
\usepackage[T1]{fontenc}
\usepackage[utf8]{inputenc}

\usepackage{amsmath,amssymb,amsthm}

\usepackage{graphicx}
\usepackage{booktabs}
\usepackage{multirow}
\usepackage{array}
\usepackage{tabularx}

\usepackage[
  font      = small,
  labelfont = bf,
  format    = plain,
  justification = justified
]{caption}
\usepackage{subcaption}

\usepackage{enumitem}
\setlist[enumerate]{leftmargin=*,itemsep=2pt,parsep=0pt}
\setlist[itemize]{leftmargin=*,itemsep=2pt,parsep=0pt}

\usepackage{xcolor}
\definecolor{linkblue}{RGB}{10,70,160}
\definecolor{citegreen}{RGB}{0,110,50}

\usepackage[
  colorlinks = true,
  linkcolor  = linkblue,
  citecolor  = citegreen,
  urlcolor   = linkblue,
  pdfborder  = {0 0 0},
  breaklinks = true
]{hyperref}
\usepackage{url}

\usepackage[round,authoryear,sort]{natbib}
\bibliographystyle{plainnat}

\usepackage{titlesec}
\titlespacing*{\section}      {0pt}{18pt plus 4pt minus 2pt}{6pt plus 2pt}
\titlespacing*{\subsection}   {0pt}{12pt plus 3pt minus 1pt}{4pt plus 1pt}
\titlespacing*{\subsubsection}{0pt}{ 9pt plus 2pt minus 1pt}{3pt plus 1pt}
\titleformat{\paragraph}[runin]{\normalfont\normalsize\bfseries}{}{0em}{}[.]

\usepackage{abstract}
\renewcommand{\abstractnamefont}{\normalfont\bfseries}
\setlength{\absleftindent}{0.5in}
\setlength{\absrightindent}{0.5in}

\usepackage{setspace}
\setstretch{1.12}

\newtheorem{hypothesis}{Hypothesis}[section]
\newtheorem{proposition}{Proposition}[section]
\newtheorem{remark}{Remark}[section]

\newcommand{\Z}[1]{\mathbb{Z}/{#1}\mathbb{Z}}
\newcommand{\Fp}{\mathbb{F}_p}
\newcommand{\W}[1]{\mathbf{W}_{#1}}
\newcommand{\dlog}{\operatorname{dlog}}

\renewcommand{\topfraction}{0.85}
\renewcommand{\textfraction}{0.10}
\renewcommand{\floatpagefraction}{0.80}
\setlength{\floatsep}{14pt plus 2pt minus 2pt}
\setlength{\textfloatsep}{16pt plus 2pt minus 4pt}
\setlength{\intextsep}{12pt plus 2pt minus 2pt}

\usepackage{fancyhdr}
\pagestyle{fancy}
\fancyhf{}
\fancyhead[L]{\small\textit{Grokking Beyond Addition}}
\fancyhead[R]{\small\textit{Mani Pal, 2026}}
\fancyfoot[C]{\small\thepage}
\renewcommand{\headrulewidth}{0.4pt}

\begin{document}

\begin{center}
  {\LARGE\bfseries Grokking Beyond Addition:\\[4pt]
   Circuit-Level Analysis of Algebraic Learning\\[4pt]
   in Transformers}

  \vspace{16pt}
  {\large Mani Pal}\\[4pt]
  {\normalsize Independent Researcher}\\[2pt]
  {\normalsize
    \href{https://github.com/groot-code24/Grokking-Circuit-Level-Analysis-of-Algebraic-Learning-in-Transformers.git}%
         {\texttt{github.com/justbytecode/Grokking-Circuit-Level-Analysis-of-Algebraic-Learning-in-Transformers}}
  }\\[2pt]
  {\normalsize March 2026}
\end{center}

\vspace{6pt}\hrule\vspace{4pt}

\begin{abstract}
Mechanistic interpretability research has shown that small transformers
trained on modular addition learn a Fourier-based \emph{clock mechanism}
and exhibit \emph{grokking}---abrupt generalisation long after training
accuracy saturates \citep{nanda2023progress}.
We extend this investigation to \textbf{eight} algebraic operations spanning
abelian fields, a composite ring, and four non-abelian groups ($S_3$, $D_5$,
$A_4$, $S_4$), training 1-layer transformers at $d_{\text{model}}=64$
with 3~seeds per operation.
Our main findings are:
(i)~all four abelian operations grokked to 100\% test accuracy within
2,000~epochs, while \emph{none} of the four non-abelian groups grokked
within the same budget despite reaching 100\% training accuracy---revealing
a clean capacity-dependent generalisation boundary;
(ii)~discrete-log re-indexing improves multiplication's Fourier concentration
by $2.14\times$ (raw $9.4\%\to$ re-indexed $20.0\%$, $g{=}3$), providing
geometric evidence for the dlog representation hypothesis;
(iii)~Peter--Weyl analysis of the non-grokked non-abelian models identifies
the correct predicted dominant irrep in all four cases ($S_3$: standard
$d{=}2$; $D_5$: $\rho_{2a}$; $A_4$: $\rho_3$, $d{=}3$; $S_4$:
standard$_3$, $d{=}3$), suggesting partial circuit formation even without
full generalisation;
(iv)~centred kernel alignment (CKA) reveals uniformly high cross-operation
embedding similarity (${\geq}\,0.80$ for all 28 pairs, mean $0.90$),
with the striking add--$S_3$ pair at $0.97$; and
(v)~the observed abelian grokking order
$\text{add}\ (1583{\pm}118)\approx\text{mul}\ (1600{\pm}187)
 < \text{sub}\ (1767{\pm}85)\approx\text{ring}\ (1767{\pm}165)$
partially confirms the complexity-delay prediction while showing subtraction
unexpectedly as slow as composite-ring addition.
Code and a checkpoint-resume Colab notebook reproduce all results in
approximately 3~hours on a free T4 GPU.
\end{abstract}

\vspace{4pt}\hrule\vspace{14pt}

\section{Introduction}
\label{sec:intro}

Grokking---delayed generalisation in small networks trained on structured,
low-data tasks---was introduced by \citet{power2022grokking} and explained
at the circuit level by \citet{nanda2023progress}, who showed that a 1-layer
transformer on modular addition discovers a Fourier clock algorithm.
\citet{chughtai2023toy} confirmed Fourier circuits for abelian groups and
began systematic non-abelian comparison.
The central open question is: \emph{what algebraic structure determines
the difficulty of grokking, and at what model capacity does generalisation
become possible?}

\paragraph{Contributions.}
\begin{enumerate}
  \item \textbf{Abelian--non-abelian grokking boundary.}
        At $d_{\text{model}}=64$, all four abelian operations grok
        (100\% rate, 3~seeds each); all four non-abelian groups
        \emph{fail} to grok (0\% rate), despite reaching 100\% training
        accuracy (\S\ref{sec:abelian}--\ref{sec:nonabelian}).
        This clean boundary is the first empirical identification of a
        capacity-dependent abelian/non-abelian grokking transition.

  \item \textbf{Partial dlog geometric evidence.}
        Multiplication embeddings show a $2.14\times$ Fourier concentration
        improvement under dlog re-indexing, consistent with the dlog
        hypothesis at this model scale (\S\ref{sec:e2}).

  \item \textbf{Pre-grokking Peter--Weyl circuit formation.}
        Non-abelian models that memorised but did not generalise already
        show the correct predicted dominant irrep, indicating that circuit
        formation begins before grokking (\S\ref{sec:nonabelian}).

  \item \textbf{CKA cross-operation analysis.}
        All 28 embedding pairs have CKA $\geq 0.80$ (mean $0.90$),
        revealing a shared representational idiom across qualitatively
        different algebraic tasks (\S\ref{sec:cka}).
        This is new relative to \citet{chughtai2023toy}.

  \item \textbf{$S_4$ Peter--Weyl analysis.}
        $S_4$ ($|G|{=}24$, 576 training pairs) is the largest group
        studied with full Peter--Weyl decomposition, removing the
        dataset-size confound of smaller groups (\S\ref{sec:nonabelian}).

  \item \textbf{Reproducibility.}
        Open checkpoint-resume Colab notebook; all results reproduced
        from checkpoints on a free T4 GPU in $\approx3$~hours.
\end{enumerate}

\section{Background}
\label{sec:background}

\subsection{Transformers and the residual stream}

We use the residual stream framing of \citet{elhage2021mathematical}.
A 1-layer transformer processes the token sequence $[a,\,b,\,\text{SEP}]$
via:
\begin{align}
  x_0 &= \W{E}[t] + \W{\mathrm{pos}}[i], \notag\\
  x_1 &= x_0 + \mathrm{Attn}(x_0), \notag\\
  x_2 &= x_1 + \mathrm{MLP}(x_1), \notag\\
  \text{logits} &= \W{U}\,x_2[\text{SEP}]. \notag
\end{align}
The prediction for the result of the algebraic operation is read from
the position of the SEP token.

\subsection{Grokking and circuit formation}

\citet{power2022grokking} introduced grokking in the context of small
algorithmic tasks. \citet{nanda2023progress} identified three qualitative
training phases---memorisation, circuit formation, and clean-up---and
reverse-engineered the Fourier clock circuit for modular addition.
\citet{varma2023explaining} provided a quantitative explanation via circuit
efficiency: weight decay penalises the large memorising solution and
eventually favours the compact generalising circuit, making grokking a
phase transition. \citet{liu2022towards} characterised grokking dynamics
via representation learning theory. \citet{chughtai2023toy} confirmed
Fourier circuits for abelian groups and performed Peter--Weyl analysis
for $S_3$, $D_5$, $A_4$.

\subsection{Fourier analysis and discrete logarithms}

For a prime $p$, the Fourier clock circuit exploits the product-to-sum
identity:
\begin{equation}
  \cos\!\Bigl(\frac{2\pi k(a+b)}{p}\Bigr)
  = \cos\!\Bigl(\frac{2\pi ka}{p}\Bigr)\cos\!\Bigl(\frac{2\pi kb}{p}\Bigr)
  - \sin\!\Bigl(\frac{2\pi ka}{p}\Bigr)\sin\!\Bigl(\frac{2\pi kb}{p}\Bigr).
  \label{eq:trig}
\end{equation}
The discrete-log hypothesis \citep{nanda2023progress} uses the isomorphism:
\begin{equation}
  \dlog_g(a \cdot b) \equiv \dlog_g(a) + \dlog_g(b) \pmod{p-1},
  \label{eq:dlog}
\end{equation}
reducing multiplication in $\mathbb{F}_p^*$ to addition in $\Z{p-1}$.

\subsection{Peter--Weyl representation theory}

For a finite group $G$, the Peter--Weyl theorem states that any function
$f: G \to \mathbb{R}^d$ decomposes into irreducible representations (irreps)
$\{\rho_k\}_{k=1}^r$ of dimensions $\{d_k\}$
\citep{diaconis1988group,serre1977linear}.
For abelian groups all irreps are one-dimensional ($d_k=1$), and the
decomposition reduces to the standard DFT.
For non-abelian groups, $d_k>1$ is possible, requiring the full
Peter--Weyl analysis; standard DFT fails to capture the multi-dimensional
irrep structure.

\section{Experimental Setup}
\label{sec:setup}

\subsection{Model architecture}

We use a 1-layer, 4-head transformer with $d_{\text{model}}=64$,
$d_{\text{head}}=16$, $d_{\text{mlp}}=256$, ReLU activations, and
layer normalisation. The input is three tokens $[a,\, b,\, \text{SEP}]$
and the prediction is taken at the SEP position.
Optimiser: AdamW, learning rate $10^{-3}$, weight decay $1.0$, full-batch.

\paragraph{Note on model scale.}
We use $d_{\text{model}}=64$ rather than the $d_{\text{model}}=128$ of
\citet{nanda2023progress}. This choice is motivated by Colab free-tier
constraints and produces qualitatively consistent results while revealing
the capacity-dependent grokking boundary. Quantitative metrics such as
Fourier concentration are lower at this scale; this is discussed explicitly
in each relevant section.

\subsection{Tasks and datasets}

Table~\ref{tab:tasks} describes all eight algebraic tasks.
All experiments use a 70/30 train/test split with a fixed random seed.
We run \textbf{3 seeds} (42, 1, 7) per operation for 3,000 epochs
(5,000 for $S_4$).

\begin{table}[t]
  \centering
  \small
  \caption{Eight algebraic tasks in order of formal complexity.
           $|\mathcal{D}|$: total dataset size.
           $d_{\max}$: maximum irrep dimension.
           $C_1$, $C_2$: formal complexity scores
           (Section~\ref{sec:delay}).}
  \label{tab:tasks}
  \renewcommand{\arraystretch}{1.25}
  \begin{tabular}{@{}clcccc@{}}
    \toprule
    \textbf{ID} & \textbf{Operation} & $|\mathcal{D}|$ &
    $d_{\max}$ & $C_1$ & $C_2$ \\
    \midrule
    E1 & $(a+b)\bmod 113$          & 12,769 & 1 & 1.1 & 1.99 \\
    E3 & $(a-b)\bmod 113$          & 12,769 & 1 & 1.1 & 1.99 \\
    E2 & $(a\times b)\bmod 113$    & 12,769 & 1 & 2.1 & 1.99 \\
    E4 & $(a+b)\bmod 100$ (ring)   & 10,000 & 1 & 3.1 & 1.99 \\
    E5 & $S_3$\quad$|G|=6$         &     36 & 2 & 4.7 & 3.25 \\
    E6 & $D_5$\quad$|G|=10$        &    100 & 2 & 5.7 & 3.34 \\
    E7 & $A_4$\quad$|G|=12$        &    144 & 3 & 6.8 & 3.75 \\
    E8 & $S_4$\quad$|G|=24$        &    576 & 3 & 7.8 & 3.80 \\
    \bottomrule
  \end{tabular}
\end{table}

\subsection{Analysis methods}

\begin{itemize}
  \item \textbf{Fourier embedding analysis.} Project $\W{E}$ onto the
        DFT basis; report L2-norm concentration in the top-5 frequency pairs.
  \item \textbf{Peter--Weyl analysis.} Compute irrep energies via
        Equation~\eqref{eq:peterwyl} from pre-computed character tables.
  \item \textbf{Discrete-log probe.} Train a linear classifier on $\W{E}$
        to predict $\dlog_g(a)$; evaluate with 5-fold cross-validation.
  \item \textbf{CKA.} Linear centred kernel alignment between embedding
        matrix pairs.
  \item \textbf{Logit attribution.} Decompose correct-answer logits into
        attention head, MLP, and embedding contributions.
  \item \textbf{Statistics.} Multi-seed results report mean~$\pm$~std
        across 3~seeds; all grok rates are exact fractions.
\end{itemize}

\section{Abelian Operations (E1--E4)}
\label{sec:abelian}

\subsection{E1: Modular Addition}
\label{sec:e1}

Figure~\ref{fig:grokking} shows the grokking dynamics for
$(a+b)\bmod 113$ across 3~seeds. Training accuracy reaches 100\% by
epoch~$\sim$200; test accuracy then plateaus near zero for over 1,000
epochs before jumping sharply to 100\% at \textbf{epoch $1583\pm118$}
(seeds: 1500, 1500, 1750).

\begin{figure}[t]
  \centering
  \includegraphics[width=0.92\textwidth]{figures/fig1b_grokking_multiseed.png}
  \caption{\textbf{E1 --- Modular addition grokking dynamics (3 seeds).}
           Solid line: mean training accuracy.
           Dashed line: mean test accuracy.
           Shaded band: $\pm1$ standard deviation across seeds.
           Red dotted line: mean grokking epoch $1583\pm118$
           (seeds: 1500, 1500, 1750).
           The $\sim$1,400-epoch gap between memorisation and
           generalisation is the defining signature of grokking.}
  \label{fig:grokking}
\end{figure}

Figure~\ref{fig:fourier} shows the Fourier embedding spectrum.
Top-5 frequencies ($k = 9, 37, 1, 29, 38$) account for \textbf{20.84\%}
of total embedding norm. This is below the $>80\%$ value reported by
\citet{nanda2023progress} at $d_{\text{model}}=128$, consistent with
the smaller embedding space distributing energy more broadly across
frequencies while still forming discernible dominant peaks. The MLP
logit attribution is consistent with a frequency-processing role.

\begin{figure}[t]
  \centering
  \includegraphics[width=0.92\textwidth]{figures/fig3_fourier_spectrum.png}
  \caption{\textbf{E1 --- Fourier embedding spectrum} ($p=113$).
           Highlighted bars (purple): top-5 frequencies
           $k\in\{9,37,1,29,38\}$.
           Top-5 concentration $= 20.84\%$.
           Dominant peaks are clearly visible above the noise floor,
           consistent with partial clock-circuit formation at
           $d_{\text{model}}=64$.}
  \label{fig:fourier}
\end{figure}

\begin{hypothesis}[Clock circuit, \citeauthor{nanda2023progress}]
  A 1-layer transformer trained on $(a+b)\bmod p$ represents each token
  $a$ as a vector containing $(\sin(2\pi ka/p),\,\cos(2\pi ka/p))$ for a
  sparse set of key frequencies $k$, and uses Eq.~\eqref{eq:trig} to
  compute the sum at the SEP position.
\end{hypothesis}

The dominant Fourier peaks and logit attribution provide \emph{partial}
support for this hypothesis at $d_{\text{model}}=64$. Full quantitative
confirmation ($>70\%$ concentration) requires $d_{\text{model}}\geq128$.

\subsection{E2: Modular Multiplication and the Discrete-Log Hypothesis}
\label{sec:e2}

\begin{hypothesis}[Discrete-log representation, \citeauthor{nanda2023progress}]
  \label{hyp:dlog}
  A transformer trained on $(a\times b)\bmod p$ encodes each token $a$
  in its embedding proportionally to $\dlog_g(a)\in\Z{p-1}$, for some
  primitive root $g$ of $\mathbb{F}_p^*$.
\end{hypothesis}

Figure~\ref{fig:dlog} presents the three-panel dlog analysis.
Applying the DFT directly over $\Z{p}$ yields a near-uniform spectrum
(concentration \textbf{9.36\%}). Re-indexing by $\dlog_g$ with $g=3$
and applying DFT over $\Z{p-1}$ reveals clear dominant peaks
(concentration \textbf{20.02\%}): a \textbf{2.14$\times$ improvement}.

\begin{figure}[t]
  \centering
  \includegraphics[width=\textwidth]{figures/fig4_dlog_analysis_panel.png}
  \caption{\textbf{E2 --- Discrete-log hypothesis test} ($g=3$, $p=113$).
           \textbf{Left:} raw Fourier spectrum, concentration $9.36\%$
           (near-uniform).
           \textbf{Centre:} dlog-reindexed Fourier spectrum, concentration
           $20.02\%$, improvement ratio $2.14\times$.
           \textbf{Right:} summary metrics. Linear probe accuracy $= 0.00$
           indicates the dlog encoding is non-linearly organised at this
           model scale; a 2-layer MLP probe is needed to recover it.}
  \label{fig:dlog}
\end{figure}

The $2.14\times$ concentration improvement constitutes \emph{geometric
evidence} for Hypothesis~\ref{hyp:dlog}: the embedding geometry aligns
with the dlog coordinate system. The linear probe accuracy of 0.00
indicates the encoding is non-linearly organised at $d_{\text{model}}=64$,
not that the encoding is absent. A non-linear (MLP) probe and the causal
wrong-root null test are implemented in the code but require
$d_{\text{model}}=128$ for interpretable results.
E2 grokked at mean epoch \textbf{$1600\pm187$}
(seeds: 1400, 1550, 1850).

\subsection{E3: Modular Subtraction (control)}

$(a-b)\bmod p$ is algebraically isomorphic to addition via token-level
negation. The circuit table (Figure~\ref{fig:circuit}) identifies a
Fourier clock representation with dominant frequencies
$k\in\{14,44,49,40\}$. Grokking at mean epoch \textbf{$1767\pm85$}
(seeds: 1800, 1650, 1850)---notably slower than addition,
discussed in Section~\ref{sec:delay}.

\subsection{E4: Ring Addition over $\mathbb{Z}/100\mathbb{Z}$}

With composite modulus $n=100$, the model produces a \emph{partial Fourier}
representation (dominant frequencies $k\in\{43,9,12,40\}$). These do not
cleanly align with the divisors of 100 at this model scale.
Grokking at mean epoch \textbf{$1767\pm165$}
(seeds: 1650, 2000, 1650).

\section{Non-abelian Groups (E5--E8): Memorisation Without Generalisation}
\label{sec:nonabelian}

\subsection{Central finding: a capacity-dependent grokking boundary}

All four non-abelian groups reached 100\% training accuracy within the
epoch budget but \emph{failed to grok}: test accuracy remained below
25\% for all seeds (Table~\ref{tab:nonabelian}).
The contrast with the abelian operations is stark:

\begin{center}
  \renewcommand{\arraystretch}{1.2}
  \begin{tabular}{lcc}
    \toprule
    & \textbf{Grok rate} & \textbf{Final train acc} \\
    \midrule
    All 4 abelian operations & \textbf{100\%} & 100\% \\
    All 4 non-abelian groups & \textbf{0\%}   & 100\% \\
    \bottomrule
  \end{tabular}
\end{center}

\begin{table}[t]
  \centering
  \small
  \caption{\textbf{Non-abelian results} at $d_{\text{model}}=64$,
           3~seeds (1~seed for $S_4$).
           All operations: train accuracy $= 100\%$; none grokked.
           Peter--Weyl dominant irrep matches the representation-theoretic
           prediction in all four cases.}
  \label{tab:nonabelian}
  \renewcommand{\arraystretch}{1.25}
  \begin{tabular}{@{}lcccc@{}}
    \toprule
    \textbf{Group} & $|G|$ & \textbf{Mean test acc} &
    \textbf{Dominant irrep} & \textbf{Evidence} \\
    \midrule
    $S_3$ &  6 &  6.1\% & standard ($d=2$)     & moderate \\
    $D_5$ & 10 &  4.4\% & $\rho_{2a}$ ($d=2$)  & moderate \\
    $A_4$ & 12 &  2.3\% & $\rho_3$ ($d=3$)     & moderate \\
    $S_4$ & 24 & 21.4\% & standard$_3$ ($d=3$) & weak     \\
    \bottomrule
  \end{tabular}
\end{table}

\subsection{Why standard Fourier analysis fails}

Raw Fourier concentration for all four groups is below 20\%, confirming
no abelian clock circuit forms. This is mathematically expected: the
abelian DFT assumes one-dimensional irreps, which do not exist for
non-abelian groups.

\subsection{Peter--Weyl analysis: circuit formation without grokking}
\label{sec:peterwyl}

Despite failing to grok, the embedding matrices show the beginnings of
the correct algebraic structure. We compute the energy of $\W{E}$ in
irrep $\rho_k$ via:
\begin{equation}
  E(\rho_k)
  \;=\; d_k \sum_{g,h\in G} K(g,h)\,\chi_k\!\left(g\cdot h^{-1}\right),
  \label{eq:peterwyl}
\end{equation}
where $K[g,h] = \W{E}[g] \cdot \W{E}[h]$ is the embedding Gram matrix
and $\chi_k$ is the character of $\rho_k$
\citep{diaconis1988group,serre1977linear}.

\paragraph{Character table note.}
For $A_4$, the complex one-dimensional irreps have character
$\chi(\text{3-cycle}) = \operatorname{Re}(e^{2\pi i/3}) = -\tfrac{1}{2}$
(not $-1$; verified against \citealt{serre1977linear}).

The dominant irrep matches the representation-theoretic prediction for all
four groups (Table~\ref{tab:nonabelian}). $S_3$, $D_5$, and $A_4$ show
moderate evidence; $S_4$ shows weak evidence, consistent with the largest
group requiring more capacity to concentrate energy in a single irrep.
Figure~\ref{fig:circuit} summarises learned circuit types for all
eight operations.

\begin{figure}[t]
  \centering
  \includegraphics[width=0.95\textwidth]{figures/fig15_circuit_table.png}
  \caption{\textbf{Learned circuit descriptions for all eight operations.}
           \textcolor{blue}{Blue rows}: abelian tasks (Fourier clock or
           dlog-then-clock).
           \textcolor{red}{Red rows}: non-abelian tasks (Peter--Weyl irrep
           projection, all memorisation-only models).
           $S_3$, $D_5$, and $A_4$ show moderate Peter--Weyl evidence
           despite never grokking, indicating partial circuit formation.}
  \label{fig:circuit}
\end{figure}

\paragraph{Interpretation.}
The Peter--Weyl circuit begins forming during the memorisation phase,
\emph{before} generalisation occurs. Non-abelian groups are in a
``stuck'' regime at $d_{\text{model}}=64$: the model discovers the correct
algebraic structure but lacks sufficient embedding dimension to complete
the generalising circuit. This is consistent with
\citet{varma2023explaining}: weight decay favours the generalising circuit,
but if that circuit requires more dimensions than the model provides,
memorisation is the stable fixed point. Increasing $d_{\text{model}}$ is
the direct test.

\section{Cross-Operation Embedding Similarity (CKA)}
\label{sec:cka}

Figure~\ref{fig:cka} shows pairwise linear CKA between the embedding
matrices of all eight trained models.

\begin{figure}[t]
  \centering
  \includegraphics[width=0.85\textwidth]{figures/fig10_cka_heatmap.png}
  \caption{\textbf{Pairwise CKA similarity of token embeddings} across
           all eight operations (28 off-diagonal pairs).
           All values $\geq 0.80$; mean $0.90$, std $0.04$.
           Highest off-diagonal: add--$S_3 = 0.97$.
           Lowest: $A_4$--$S_4 = 0.80$.
           The uniformly high similarity suggests a shared representational
           substrate across abelian and non-abelian tasks at this scale.}
  \label{fig:cka}
\end{figure}

All 28 pairwise CKA values lie in $[0.80,\, 0.97]$ (mean $0.90$,
std $0.04$). Three observations follow.

\paragraph{Add--$S_3 = 0.97$.}
The highest off-diagonal value pairs modular addition (grokked, epoch 1583)
with $S_3$ (memorised only, test accuracy 6.1\%). A fully generalised
abelian model and a memorised non-abelian model share nearly identical
embedding geometry. This suggests the early-stage embedding structure is
driven by factors shared across all tasks---most likely positional embeddings,
token frequency, and the cross-entropy loss---rather than task-specific
algebraic circuits.

\paragraph{Uniformly high across abelian--non-abelian pairs.}
Mean CKA for abelian--non-abelian pairs is 0.92; for non-abelian pairs
it is 0.87. The embedding geometry is largely determined by shared
features at $d_{\text{model}}=64$.

\paragraph{$A_4$--$S_4 = 0.80$ is the global minimum.}
The two largest non-abelian groups diverge most from each other and from
abelian operations, consistent with their partial but incomplete circuit
formation.

\begin{remark}
The uniformly high CKA motivates a capacity-scaling study. If the high
similarity reflects capacity limitation rather than genuine representational
universality, CKA values should decrease and differentiate at larger
$d_{\text{model}}$. This is a testable prediction.
\end{remark}

\section{Abelian Grokking Delays and Complexity}
\label{sec:delay}

\subsection{Observed delays}

Figure~\ref{fig:cdf} and Table~\ref{tab:delay} summarise grokking CDFs
for all four abelian operations across 3~seeds each.

\begin{figure}[t]
  \centering
  \includegraphics[width=0.92\textwidth]{figures/fig16_grokking_cdf.png}
  \caption{\textbf{Grokking epoch CDF for abelian operations} (3 seeds
           each, all 100\% grok rate).
           Addition (purple) and multiplication (gold) are nearly
           indistinguishable in timing (mean 1583 vs.\ 1600).
           Subtraction (green) and ring addition (orange) tie at
           mean epoch 1767.}
  \label{fig:cdf}
\end{figure}

\begin{table}[t]
  \centering
  \small
  \caption{\textbf{Abelian grokking delays.}
           All operations: 100\% grok rate, 3~seeds.
           $C_1$/$C_2$: formal complexity scores.
           $\sigma$: standard deviation across seeds.}
  \label{tab:delay}
  \renewcommand{\arraystretch}{1.25}
  \begin{tabular}{@{}lccccc@{}}
    \toprule
    \textbf{Operation} & $C_1$ & $C_2$ & \textbf{Mean} & $\sigma$ &
    \textbf{Seeds} \\
    \midrule
    Addition       & 1.1 & 1.99 & 1583 & 118 & 1500,\ 1500,\ 1750 \\
    Multiplication & 2.1 & 1.99 & 1600 & 187 & 1400,\ 1550,\ 1850 \\
    Subtraction    & 1.1 & 1.99 & 1767 &  85 & 1800,\ 1650,\ 1850 \\
    Ring addition  & 3.1 & 1.99 & 1767 & 165 & 1650,\ 2000,\ 1650 \\
    \bottomrule
  \end{tabular}
\end{table}

\subsection{Formal complexity scores}

\paragraph{$C_1$ (ordinal rank).}
\begin{equation}
  C_1(\mathcal{A}) \;=\;
  \mathrm{rank}(\mathcal{A})
  + \frac{\max_k d_k}{10}
  + \frac{1}{2}\,\mathbf{1}[\text{non-abelian}].
  \label{eq:c1}
\end{equation}

\paragraph{$C_2$ (theory-grounded).}
\begin{equation}
  C_2(G) \;=\;
  \log_2\!\bigl(\max_k d_k + 1\bigr)
  + \left(1 - \frac{1}{|\widehat{G}|}\right)
  + \mathbf{1}[\text{non-abelian}],
  \label{eq:c2}
\end{equation}
where $|\widehat{G}|$ is the number of irreducible representations.
The three terms capture: (1)~the information-theoretic cost of representing
the hardest irrep; (2)~irrep diversity; and (3)~the abelian-to-non-abelian
transition. Both scores are computable from character tables alone, without
reference to training outcomes.

\subsection{Comparison with predictions}

The formal scores predict:
$\text{add}\approx\text{sub}<\text{mul}<\text{ring}<\text{non-abelian}$.
The observed order is:
$\text{add}\ (1583)\approx\text{mul}\ (1600)<\text{sub}\ (1767)\approx
\text{ring}\ (1767)\ll\text{non-abelian (did not grok)}$.

\noindent\textbf{Three predictions confirmed:}
\begin{itemize}
  \item Ring addition is the hardest abelian task:
        $1767 > 1583$ and $1767 > 1600$. \checkmark
  \item All non-abelian tasks are harder than all abelian tasks:
        0\% vs.\ 100\% grok rate. \checkmark
  \item The ordering ring $>$ add is consistent across all seeds. \checkmark
\end{itemize}

\noindent\textbf{Two deviations observed:}
\begin{itemize}
  \item \textbf{Add $\approx$ mul} (predicted: mul $>$ add).
        The 17-epoch mean difference is well within overlapping error bars
        ($\sigma_{\text{add}}=118$, $\sigma_{\text{mul}}=187$) and is not
        statistically resolvable at 3~seeds.

  \item \textbf{Sub $\approx$ ring} (predicted: sub $\approx$ add $\ll$ ring).
        Subtraction grokked 184~epochs later than addition despite being
        algebraically isomorphic via negation. We conjecture this reflects
        the cost of discovering the token-level negation map from data,
        adding a seed-independent overhead.
\end{itemize}

\begin{proposition}[Partial complexity--delay law]
  \label{prop:delay}
  At $d_{\text{model}}=64$ with 3~seeds, the formal scores $C_1$/$C_2$
  correctly predict: (1)~ring addition as the hardest abelian operation;
  and (2)~all non-abelian operations as strictly harder than all abelian
  operations. The add~$\approx$~mul and sub~$\approx$~ring ties are
  statistically unresolved and require at least 5~seeds or
  $d_{\text{model}}\geq128$.
\end{proposition}

\section{Related Work}
\label{sec:related}

\citet{power2022grokking} introduced grokking. \citet{nanda2023progress}
reverse-engineered the clock circuit and proposed the dlog hypothesis for
multiplication. \citet{chughtai2023toy} confirmed Fourier circuits for
abelian groups and conducted Peter--Weyl analysis for $S_3$, $D_5$, $A_4$.
\citet{liu2022towards} characterised grokking dynamics via representation
learning theory. \citet{varma2023explaining} explained grokking via circuit
efficiency and weight decay. \citet{barak2022hidden} found hidden-progress
phenomena in learning parities. \citet{elhage2021mathematical} introduced
the residual stream decomposition. \citet{conmy2023towards} developed
automated circuit discovery. \citet{diaconis1988group} and
\citet{serre1977linear} provide the mathematical foundations for Peter--Weyl
and character-table computations. \citet{cohen2016group} showed
group-equivariant CNNs encode group structure by architectural design; we
show transformers discover it without any such bias.

\paragraph{Novelty relative to prior work.}
\citet{chughtai2023toy} covers $S_3$, $D_5$, $A_4$ but not $S_4$.
To our knowledge, this paper is the first to provide:
(a)~systematic Peter--Weyl analysis of $S_4$ ($|G|{=}24$, 576 training pairs);
(b)~CKA cross-operation similarity measurement spanning both abelian and
non-abelian tasks; and
(c)~empirical identification of the capacity-dependent abelian/non-abelian
grokking boundary at $d_{\text{model}}=64$.

\section{Conclusion}
\label{sec:conclusion}

\paragraph{Established results.}
All four abelian operations grok reliably at $d_{\text{model}}=64$
(100\% rate, all within 2,000 epochs); all four non-abelian groups
memorise but fail to grok (0\% rate, 100\% train accuracy).
This abelian/non-abelian grokking boundary, consistent with the
circuit-efficiency theory of \citet{varma2023explaining}, is the central
empirical finding of this paper. Non-abelian models develop the correct
Peter--Weyl circuit signature even without generalising, demonstrating
that circuit formation is separable from successful grokking. CKA reveals
uniformly high cross-operation embedding similarity ($\geq0.80$,
mean $0.90$), with the striking add--$S_3$ pair at $0.97$. Discrete-log
re-indexing provides a $2.14\times$ Fourier concentration improvement for
multiplication, supporting the dlog hypothesis at the geometric level.

\paragraph{Open questions.}
\begin{enumerate}
  \item Do non-abelian groups grok at $d_{\text{model}}=128$ or $256$?
        The partial Peter--Weyl circuit formation observed here suggests yes.
  \item Is the add~$\approx$~mul tie real or a 3-seed artefact?
        Requires at least 5~seeds.
  \item Does the high CKA reflect genuine representational universality or
        purely capacity limitation?
  \item What is the causal dlog accuracy and null-test $p$-value at
        $d_{\text{model}}=128$?
\end{enumerate}

\paragraph{Summary thesis.}
At capacity $d_{\text{model}}=64$, 1-layer transformers readily grok
abelian algebraic tasks but encounter a hard generalisation boundary for
non-abelian tasks. The boundary is consistent with representation theory:
abelian groups have one-dimensional irreps that fit naturally in a small
embedding space; non-abelian groups require higher-dimensional irreps that
exceed this capacity. Scaling $d_{\text{model}}$ is the direct,
falsifiable test of this interpretation.

\section{Limitations}
\label{sec:limitations}

\begin{enumerate}
  \item \textbf{Model scale.} $d_{\text{model}}=64$ is smaller than
        \citet{nanda2023progress}. Fourier concentration (20.84\%) and
        dlog linear probe accuracy (0.00) are capacity-limited; quantitative
        comparison with prior results requires $d_{\text{model}}=128$.

  \item \textbf{Seed count.} 3~seeds is sufficient for binary
        grok/no-grok classification but borderline for ordering statistics.
        The add~$\approx$~mul tie requires at least 5~seeds.

  \item \textbf{Non-abelian non-grokking.} All four non-abelian groups
        failed to grok. We cannot characterise their final generalising
        circuit because no generalising model was obtained at this scale.

  \item \textbf{Causal dlog verification.} The wrong-root null test is
        implemented but numeric accuracy and $p$-value results await the
        $d_{\text{model}}=128$ run.

  \item \textbf{Ablations not executed.} Dataset-size, training-fraction,
        and weight-decay ablations are implemented in the codebase but were
        not run in this experiment.

  \item \textbf{Inductive claims.} All conclusions are empirical from
        8~operations; no formal proofs are provided.
\end{enumerate}

\section*{Acknowledgements}

We thank Neel Nanda for TransformerLens and the mechanistic interpretability
tutorials that form the foundation of this work.

\bibliography{references}

\end{document}